% Part: second-order-logic
% Chapter: syntax-and-semantics
% Section: introduction

\documentclass[../../../include/open-logic-section]{subfiles}

\begin{document}

\olfileid{sol}{syn}{int}

\olsection{مقدمه}

در منطق مرتبهٔ اول، نمادهای غیرمنطقیِ یک زبان معین، یعنی
!!{constant}s، !!{function}s و !!{predicate}s آن، را با نمادهای منطقی
ترکیب می‌کنیم تا مطالبی دربارهٔ !!{structure}s مرتبهٔ اول بیان کنیم.
این کار با استفاده از مفهوم ارضا انجام می‌شود که میان یک
!!a{structure} چون~$\Struct{M}$، همراه با یک تخصیص متغیر چون~$s$، و
!!a{formula} چون~$!A$ رابطه برقرار می‌کند: $\Sat{M}{!A}[s]$ برقرار
است اگر و تنها اگر آنچه~$!A$، با تفسیر !!{constant}s، !!{function}s و
!!{predicate}s خود مطابق~$\Struct{M}$ و تفسیر متغیرهای آزادش مطابق~$s$،
بیان می‌کند صادق باشد. تفسیر !!{identity} یعنی~$\eq$، در تعریف
$\Sat{M}{!A}[s]$ گنجانده شده است؛ تفسیر~$\lforall$ و~$\lexists$ نیز چنین است.
اولی همواره به‌صورت رابطهٔ همانی بر !!{domain}~$\Domain{M}$ ساختار
تفسیر می‌شود و سورها همواره بر سراسر !!{domain} دامنه دارند. اما نکتهٔ
اساسی این است که کمیت‌گذاری تنها بر عناصر !!{domain} مجاز است؛ ازاین‌رو
فقط !!{variable}s شیئی می‌توانند پس از یک سور بیایند.

در منطق مرتبهٔ دوم، هم زبان و هم تعریف ارضا گسترش می‌یابند تا متغیرهای
تابعی و محمولیِ آزاد و مقید و کمیت‌گذاری بر آن‌ها را در بر گیرند. نسبت
این متغیرها با !!{function}s و !!{predicate}s همانند نسبت متغیرهای شیئی
با !!{constant}s است. آن‌ها در تشکیل ترم‌ها و !!{formula}s منطق مرتبهٔ
دوم همان نقش را ایفا می‌کنند و کمیت‌گذاری بر آن‌ها نیز به شیوه‌ای مشابه
انجام می‌شود. در معناشناسی \emph{استاندارد}، سورهای مرتبهٔ دوم بر همهٔ
اشیای ممکن از نوع مناسب دامنه دارند (برای متغیرهای تابعی، توابع
$n$-موضعی از~$\Domain{M}^n$ به~$\Domain{M}$، و برای متغیرهای محمولی،
روابط $n$-موضعی). برای نمونه، هرچند
$\lforall[\Obj{v_0}][(\Obj{P^1_0}(\Obj{v_0}) \lor \lnot
  \Obj{P^1_0}(\Obj{v_0}))]$ هم در منطق مرتبهٔ اول و هم در منطق مرتبهٔ
دوم فرمول است، در دومی می‌توانیم
$\lforall[\Obj{V^1_0}][\lforall[\Obj{v_0}][(\Obj{V^1_0}(\Obj{v_0}) \lor
    \lnot \Obj{V^1_0}(\Obj{v_0}))]]$ و
$\lexists[\Obj{V^1_0}][\lforall[\Obj{v_0}][(\Obj{V^1_0}(\Obj{v_0}) \lor
    \lnot \Obj{V^1_0}(\Obj{v_0}))]]$ را نیز در نظر بگیریم. چون این
عبارت‌ها متغیر آزاد ندارند، !!{sentence}s منطق مرتبهٔ دوم‌اند. در اینجا
$\Obj{V^1_0}$ یک متغیر محمولیِ $1$-موضعیِ مرتبهٔ دوم است. تفسیرهای
مجاز~$\Obj{V^1_0}$ همان چیزهایی‌اند که می‌توانیم به یک !!{predicate}
$1$-موضعی چون~$\Obj{P^1_0}$ نسبت دهیم، یعنی زیرمجموعه‌های
$\Domain{M}$. کمیت‌گذاری بر آن‌ها بدین معناست که
$\lforall[\Obj{v_0}][(\Obj{V^1_0}(\Obj{v_0}) \lor \lnot
  \Obj{V^1_0}(\Obj{v_0}))]$ برای همهٔ شیوه‌های تخصیص یک زیرمجموعه از
$\Domain{M}$ به‌عنوان ارزش~$\Obj{V^1_0}$، یا دست‌کم برای یکی از آن‌ها،
برقرار است. چون هر مجموعه یا شیئی معین را در بر می‌گیرد یا در بر
نمی‌گیرد، هر دو جمله در هر !!{structure}ی صادق‌اند.
\end{document}
